\subsection{SUBSTITUTIONS IN THE HAUSDORFF SERIES}

As K is a field containing Q, the Hausdorff series can be considered as a Lie formal power series with coefficients in K. Therefore, if g is a complete Hausdorff filtered Lie algebra with \(g = \bigcup_{\alpha > 0} g_{\alpha}\) , then, for a, b in g, a and b can be substituted for U and V in H (cf.~no. 3 and §2, no. 5, Remark).

In particular, let A be a complete Hausdorff filtered unital associative algebra. We write \(m = \bigcup_{\alpha > 0} A_{\alpha}\) and \(m_{\alpha} = A_{\alpha} \cap m\) for \(\alpha \in R\) ; hence \(m_{\alpha} = A_{\alpha}\) for \(\alpha > 0\) and \(m_{\alpha} = m\) for \(\alpha \leqslant 0\) . With the bracket \([a, b] = ab - ba\) , m is a complete Hausdorff filtered Lie algebra, to which the above can be applied. With this notation, we have the following result which completes Proposition 1 of no. 1.

PROPOSITION 3. If \(a \in \mathfrak{m}\) , \(b \in \mathfrak{m}\) , then \(\exp \mathrm{H}(a, b) = \exp a \cdot \exp b\) .

Let \(a, b\) be in \(\mathfrak{m}\) ; there exists \(\alpha > 0\) such that \(a \in \mathrm{A}_{\alpha}\) and \(b \in \mathrm{A}_{\alpha}\) . Then there exists a continuous homomorphism \(\theta\) of the Magnus algebra \(\hat{\mathbf{A}}(\{\mathbf{U},\mathbf{V}\})\) into A mapping U to \(a\) and V to \(b\) ( \(\S 5\) , no. 1, Proposition 1).

The restriction of \(\theta\) to \(\hat{\mathbf{L}}(\{\mathrm{U},\mathrm{V}\})\) is a continuous homomorphism of Lie algebras of \(\mathbf{L}(\{\mathrm{U},\mathrm{V}\})\) into \(\mathfrak{m}\) which maps \(\mathbf{U}\) (resp. V) to \(a\) (resp. \(b\) ). By formula (6) of no. 3, therefore \(\theta (\mathrm{H}) = \mathrm{H}(a,b)\) . It then suffices to apply the continuous homomorphism \(\theta\) to the two sides of the relation

\[
\exp \mathrm{H} (\mathrm{U}, \mathrm{V}) = \exp \mathrm{U}. \exp \mathrm{V}
\]

taking account of Remark 3 of no. 1.

Remark 1. If \(a\) and \(b\) commute, then \(\mathrm{H}_{r,s}(a,b) = 0\) for \(r + s \geqslant 2\) , for every homogeneous Lie polynomial of degree \(\geqslant 2\) is zero at \((a,b)\) . Then \(\mathrm{H}(a,b) = a + b\) and Proposition 3 recovers the formula

\[
\exp (a + b) = \exp a. \exp b.
\]

PROPOSITION 4. Let \(\mathfrak{g}\) be a complete Hausdorff filtered Lie algebra such that \(\mathfrak{g} = \bigcup_{\alpha >0} \mathfrak{g}_{\alpha}\) . The mapping \((a, b) \mapsto \mathrm{H}(a, b)\) is a group law on \(\mathfrak{g}\) compatible with the topology on \(\mathfrak{g}\) under which \(0\) is the identity element and \(-a\) is the inverse of \(a\) for all \(a \in \mathfrak{g}\) .

The mapping \((a, b) \mapsto \mathrm{H}(a, b)\) of \(\mathfrak{g} \times \mathfrak{g}\) into \(\mathfrak{g}\) is continuous (no. 3); as the mapping \(a \mapsto -a\) is obviously continuous, it suffices to prove the relations

\begin{enumerate}
\def\labelenumi{(\arabic{enumi})}
\setcounter{enumi}{17}
\tightlist
\item
\end{enumerate}

\[
\mathrm{H} (\mathrm{H} (a, b), c) = \mathrm{H} (a, \mathrm{H} (b, c))\tag{18}
\]

\[
\mathrm{H} (a, - a) = 0\tag{19}
\]

\[
\mathrm{H} (a, 0) = \mathrm{H} (0, a) = a\tag{20}
\]

for \(a, b, c\) in \(\mathfrak{g}\) . By formula (7) of no. 3, it suffices to prove these formulae when \(a, b, c\) are three indeterminates and \(\mathfrak{g} = \hat{\mathbf{L}}\langle \{a, b, c\}\rangle\) . Now the restriction of the exponential mapping to \(\hat{\mathbf{L}}\langle \{a, b, c\}\rangle\) is an injection into the Magnus algebra \(\hat{\mathbf{A}}\langle \{a, b, c\}\rangle\) and by Proposition 3:

\[
\exp \mathrm{H} (\mathrm{H} (a, b), c) = \exp \mathrm{H} (a, b). \exp c = \exp a. \exp b. \exp c
\]

\[
\exp \mathrm{H} (a, \mathrm{H} (b, c)) = \exp a. \exp \mathrm{H} (b, c) = \exp a. \exp b. \exp c
\]

\[
\exp \mathrm{H} (a, - a) = \exp a. \exp (- a) = \exp (a - a) = \exp 0
\]

\[
\exp \mathrm{H} (a, 0) = \exp a. \exp 0 = \exp a
\]

\[
\exp \mathrm{H} (0, a) = \exp 0. \exp a = \exp a.
\]

This establishes relations (18) to (20).

Remarks. (2) Take g to be the Lie algebra \(\hat{\mathbf{L}}(\mathbf{X})\) . The group law introduced in the above proposition coincides with the law defined in no. 2. In other words,

\[
a \uplus b = \mathrm{H} (a, b) \quad \text { for   } a, b \text {   in   } \hat {\mathrm{L}} (\mathrm{X});\tag{21}
\]

thus the Hausdorff group law is given by the Hausdorff series.

\begin{enumerate}
\def\labelenumi{(\arabic{enumi})}
\setcounter{enumi}{2}
\tightlist
\item
  Let \(\mathfrak{g}\) be a Lie algebra with the integral filtration \((\mathcal{C}^m\mathfrak{g})\) defined by the lower central series. Suppose that there exists \(m \geqslant 1\) such that \(\mathcal{C}^m\mathfrak{g} = \{0\}\) . With the topology derived from the filtration \((\mathcal{C}^m\mathfrak{g})_{n \geqslant 1}\) , the Lie algebra \(\mathfrak{g}\) is Hausdorff, complete and even discrete. Then \(\mathrm{P}(a_1, \ldots, a_r) = 0\) for \(a_1, \ldots, a_r\) in \(\mathfrak{g}\) and for every homogeneous Lie polynomial \(\mathrm{P}\) of degree \(\geqslant m\) ; in particular, \(\mathrm{H}_{r,s}(a, b) = 0\) for \(r + s \geqslant m\) and the series \(\mathrm{H}(a, b) = \sum_{r,s} \mathrm{H}_{r,s}(a, b)\) has only a finite number of non-zero terms. The group law \((a, b) \mapsto \mathrm{H}(a, b)\) on \(\mathfrak{g}\) is then a polynomial mapping (§ 2, no. 4).
\end{enumerate}

PROPOSITION 5. Let \(\mathbf{K}_{r,s}\) be the component of \(\mathrm{H}(\mathrm{U} + \mathrm{V}, -\mathrm{U})\) of multidegree \((r,s)\) . Then

\[
\mathrm{K} _ {n, 1} (\mathrm{U}, \mathrm{V}) = \frac {1}{(n + 1) !} (\operatorname{ad} \mathrm{U}) ^ {n} (\mathrm{V}) \quad for n \geqslant 0.
\]

We write \(\mathrm{K}(\mathrm{U},\mathrm{V}) = \mathrm{H}(\mathrm{U} + \mathrm{V}, - \mathrm{U}),\mathrm{K}_1(\mathrm{U},\mathrm{V}) = \sum_{n\geqslant 0}\mathrm{K}_{n,1}(\mathrm{U},\mathrm{V}).\) We denote by L (resp. R) left (resp. right) multiplication by U on \(\hat{\mathbf{A}} (\{\mathbf{U},\mathbf{V}\})\)

We can write

\[
\begin{array}{r l} e ^ {\mathrm{U}} \mathrm{V} e ^ {- \mathrm{U}} & = \sum_ {p, q} \frac {\mathrm{U} ^ {p}}{p !} \mathrm{V} \frac {(- \mathrm{U}) ^ {q}}{q !} \\ & = \sum_ {n \geqslant 0} \frac {1}{n !} \left(\sum_ {p + q = n} \frac {n !}{p ! q !} (\mathrm{L} ^ {p} (- \mathrm{R}) ^ {q}). \mathrm{V}\right) \\ & = \sum_ {n \geqslant 0} \frac {1}{n !} (\mathrm{L} - \mathrm{R}) ^ {n}. \mathrm{V} \end{array}
\]

and therefore

\[
e ^ {\mathrm{U}} \mathrm{V} e ^ {- \mathrm{U}} = \sum_ {n \geqslant 0} \frac {1}{n !} (\mathrm{adU}) ^ {n} \mathrm{V}.\tag{22}
\]

We now calculate modulo the ideal \(\sum_{m\geqslant 0}\sum_{n\geqslant 2}\mathrm{A}^{m,n}(\{\mathrm{U},\mathrm{V}\})\) of \(\mathrm{A}(\{\mathrm{U},\mathrm{V}\})\) . For all \(n\geqslant 1\) ,

\[
(\mathrm{U} + \mathrm{V}) ^ {n} \equiv \mathrm{U} ^ {n} + \sum_ {i = 1} ^ {n - 1} \mathrm{U} ^ {i} \mathrm{V} \mathrm{U} ^ {n - 1 - i}
\]

whence

\[
\begin{array}{r l} (\mathrm{adU}) (\mathrm{U} + \mathrm{V}) ^ {n} & \equiv ((\mathrm{L} - \mathrm{R}) \sum_ {i = 1} ^ {n - 1} \mathrm{L} ^ {i} \mathrm{R} ^ {n - i}). \mathrm{V} \\ & \equiv (\mathrm{L} ^ {n} - \mathrm{R} ^ {n}). \mathrm{V} \\ & \equiv \mathrm{U} ^ {n} \mathrm{V} - \mathrm{VU} ^ {n}. \end{array}
\]

Therefore

\[
(\mathrm{ad} \mathrm{U}). \mathrm{e} ^ {\mathrm{U} + \mathrm{V}} \equiv e ^ {\mathrm{U}} \mathrm{V} - \mathrm{V} e ^ {\mathrm{U}}\tag{23}
\]

summing over n.

On the other hand, \(\mathrm{K}_1(\mathrm{U},\mathrm{V})\equiv \mathrm{K}(\mathrm{U},\mathrm{V})\) and \(e^{\mathrm{K}_1(\mathrm{U},\mathrm{V})}\equiv 1 + \mathrm{K}_1(\mathrm{U},\mathrm{V})\) and hence

\[
\mathrm{K} _ {1} \equiv e ^ {\mathrm{K}} - 1 \equiv e ^ {\mathrm{U} + \mathrm{V}} e ^ {- \mathrm{U}} - 1
\]

by Proposition 3. We deduce that

\[
\begin{array}{r l r} (\mathrm{adU}) \mathrm{K} _ {1} & \equiv \mathrm{Ue} ^ {\mathrm{U} + \mathrm{V}} e ^ {- \mathrm{U}} - e ^ {\mathrm{U} + \mathrm{V}} e ^ {- \mathrm{U}} \mathrm{U} \equiv (\mathrm{Ue} ^ {\mathrm{U} + \mathrm{V}} - e ^ {\mathrm{U} + \mathrm{V}} \mathrm{U}) e ^ {- \mathrm{U}} \\ & \equiv (e ^ {\mathrm{U}} \mathrm{V} - \mathrm{Ve} ^ {\mathrm{U}}) e ^ {- \mathrm{U}} & \text {by (23)} \\ & \equiv e ^ {\mathrm{U}} \mathrm{Ve} ^ {- \mathrm{U}} - \mathrm{V} \\ & \equiv \sum_ {n \geqslant 1} \frac {1}{n !} (\mathrm{adU}) ^ {n} \mathrm{V} & \text {by (22)} \\ & \equiv (\mathrm{adU}) \Bigl (\sum_ {n \geqslant 0} \frac {(\mathrm{adU}) ^ {n}}{(n + 1) !} \mathrm{V} \Bigr). \end{array}
\]

It then suffices to apply the Remark of § 2, no. 11.

